\chapter{Objetivos}

\section{Objetivo general}
Desarrollar un sistema de apoyo a la toma de decisiones basado en visión computacional y un esquema de valoración económica para la clasificación de residuos de aparatos eléctricos y electrónicos (RAEE), ejecutable en hardware de cómputo de bajo costo, orientado a mejorar la precisión, los tiempos de clasificación y la viabilidad del proceso frente a métodos manuales.

\subsection{Objetivos específicos}
\begin{itemize}
    \item Diseñar un modelo de visión computacional basado en técnicas de transfer learning para la identificación y clasificación de RAEE en imágenes.
    \item Definir un esquema de valoración económica basado en la clasificación de RAEE y su potencial de reciclabilidad, para su categorización en niveles de valorización.
    \item Implementar y optimizar el modelo de visión computacional mediante técnicas de cuantización y conversión a formatos de inferencia ligeros, para su ejecución en una plataforma de cómputo de bajo costo (Raspberry Pi 5, con y sin acelerador de inferencia neuronal), evaluando el compromiso entre precisión, latencia y consumo de recursos.
    \item Integrar el modelo de visión computacional y el esquema de valoración en un prototipo de aplicación que genere recomendaciones de reciclaje o disposición del residuo.
    \item Evaluar el rendimiento del sistema mediante métricas de precisión y tiempo de clasificación, con el modelo ejecutándose sobre el hardware de bajo costo, en comparación con un proceso manual en un entorno controlado.
\end{itemize}
